\documentclass[a4paper,11pt]{article}
\usepackage[margin=23mm]{geometry}
\usepackage{amsmath,amssymb,bm,booktabs}
\usepackage[hidelinks]{hyperref}
\newcommand{\ii}{\mathrm{i}}
\newcommand{\Tr}{\operatorname{Tr}}
\newcommand{\ket}[1]{\lvert #1\rangle}
\newcommand{\bra}[1]{\langle #1\rvert}
\newcommand{\Diss}{\mathcal D}
\newcommand{\Ham}{\mathcal H}
\newcommand{\id}{\mathcal I}
\setlength{\parskip}{0.25em}
\setlength{\parindent}{1em}
\allowdisplaybreaks

% Working Methods draft, 2026-10-04.
% The body between BEGIN/END METHODS can be copied into the manuscript.
% Current scope: single-mode first-sideband model, rectangular-pulse
% amplitude feedback, Fock-resolved diagnostics, and Pauli-channel distances.
% Source mapping and checks still needed are in methods_draft_notes.md.

\begin{document}
\setcounter{section}{1}
% BEGIN METHODS
\section{Methods}
\label{sec:methods}

We characterize the errors of a two-ion M{\o}lmer--S{\o}rensen (MS) gate
by reconstructing its reduced two-qubit channel from master-equation
simulations. The ideal gate is factored out before analyzing the process
matrix and its logarithmic error generator. We then calibrate the drive
amplitude, repeat the channel reconstruction, and compare the remaining
noise with its Pauli-twirled approximation. A separate calculation resolved
by the initial motional Fock state is used to identify the contribution of
thermal motion to the error structure.

\subsection{Effective spin--motion Hamiltonian}
\label{subsec:hamiltonian}

We consider two qubits coupled to a single collective motional mode.
Let $a$ and $a^\dagger$ be its ladder operators, $N=a^\dagger a$, and
$S_x=X\otimes I+I\otimes X$. With a fixed spin phase, the effective
interaction-picture Hamiltonian is
\begin{equation}
 H(t)=\hbar A(t)S_x\otimes
 \left[B e^{-\ii\Phi(t)}+B^\dagger e^{\ii\Phi(t)}\right],
 \qquad \Phi(t)=\int_0^t\delta(t')\,dt',
 \label{eq:methods_hamiltonian}
\end{equation}
where $A(t)$ is the effective sideband coupling in angular-frequency units
and $\delta(t)$ is the detuning from the motional sideband. The leading
Lamb--Dicke factor is already absorbed into $A(t)$. To retain the leading
relative correction to the first-sideband coupling, we use
\begin{equation}
 B=a-\frac{\eta^2}{2}(N+1)a
   =a\left(1-\frac{\eta^2}{2}N\right),
 \qquad
 B\ket{n}=\sqrt n\left(1-\frac{\eta^2 n}{2}\right)\ket{n-1},
 \label{eq:methods_B}
\end{equation}
where $\eta$ is the Lamb--Dicke parameter. Thus the model retains an
$O(\eta^2)$ correction relative to the linear force. It does not include
the full optical interaction to all orders or additional motional
sidebands. The leading-order reference model is obtained by setting
$B=a$ at the same effective coupling.

For the baseline rectangular pulse, $A(t)=A_0$ and $\delta(t)=\delta>0$,
with $T=2\pi/\delta$. In the leading-order model and in the absence of
dissipation, the motional trajectory closes at $T$, giving, up to a global
phase,
\begin{equation}
 U_{\mathrm{LD}}(T)=\exp(\ii\phi_{\mathrm{LD}}XX),
 \qquad \phi_{\mathrm{LD}}=4\pi(A_0/\delta)^2,
 \qquad XX=X\otimes X.
\end{equation}
We choose $A_0/\delta=1/4$ and define the target operation as
$U_\star=\exp(\ii\phi_\star XX)$ with $\phi_\star=\pi/4$.
All rotation-angle coefficients below use this convention, without a
factor of $1/2$ in the exponent.

\subsection{Initial state and dissipative dynamics}
\label{subsec:master_equation}

The spin and motion are initially uncorrelated. For each input spin state
$\rho$, we initialize the motion in a thermal state,
\begin{equation}
 \varrho(0)=\rho\otimes\rho_{\mathrm{th}}(\bar n),\qquad
 \rho_{\mathrm{th}}(\bar n)=\sum_{n=0}^{\infty}p_n\ket n\bra n,
 \qquad p_n=\frac{\bar n^n}{(1+\bar n)^{n+1}}.
 \label{eq:methods_thermal}
\end{equation}
The $\bar n=0$ case is understood as the motional ground state.
The joint density operator evolves according to
\begin{align}
 \dot\varrho={}&-\frac{\ii}{\hbar}[H(t),\varrho]
 +\Gamma_{\mathrm h}\bigl(\Diss[a]+\Diss[a^\dagger]\bigr)(\varrho)
 +\Gamma_{\mathrm m}\Diss[N](\varrho) \nonumber\\
 &+\sum_{j=1}^2\left[
 \left(\frac{\Gamma_{\mathrm s}}{2}+\Gamma_{\mathrm{Ray}}\right)
       \Diss[Z_j]
 +\frac{\Gamma_{\mathrm{Ram}}}{2}
       \bigl(\Diss[\sigma_j^+]+\Diss[\sigma_j^-]\bigr)
 \right](\varrho),
 \label{eq:methods_master}
\end{align}
where
\begin{equation}
 \Diss[L](\varrho)=L\varrho L^\dagger
   -\frac12\{L^\dagger L,\varrho\}.
\end{equation}
Identity operators on unaffected subsystems are implicit.
Here $\Gamma_{\mathrm h}$ describes motional diffusion, with
$d\langle N\rangle/dt=\Gamma_{\mathrm h}$ for this dissipator alone in
an untruncated oscillator; $\Gamma_{\mathrm m}$ describes motional
dephasing; and $\Gamma_{\mathrm s}=1/T_2^*$ describes independent spin
dephasing. Rayleigh and Raman scattering are represented by effective
two-level dephasing and symmetric spin-flip terms, respectively. Leakage
out of the qubit subspace is not included.

The initial occupation $\bar n$ and the dynamical noise rates are
independent model parameters. In particular, an initially thermal mode
does not imply that the heating terms describe equilibration with a bath
at the same temperature. The baseline comparison has no additional
quasi-static fluctuations of the drive amplitude, detuning, or spin phase.

\subsection{Channel reconstruction and removal of the ideal gate}
\label{subsec:channel_reconstruction}

We reconstruct the physical two-qubit channel
$\mathcal G(\rho)=\Tr_{\mathrm m}[\varrho(T)]$
using numerical quantum process tomography~\cite{QPT}. The input set
contains the 16 tensor products of the four single-qubit states
$\ket0\bra0$, $\ket1\bra1$, $\ket+\bra+$, and
$\ket{+\ii}\bra{+\ii}$, where
$\ket+=(\ket0+\ket1)/\sqrt2$ and
$\ket{+\ii}=(\ket0+\ii\ket1)/\sqrt2$.
The same initial thermal state and dynamical parameters are used for
every input. If $M_{\mathrm{in}}$ and $M_{\mathrm{out}}$ contain the
column-vectorized input and output density matrices, respectively, the
superoperator is $S_{\mathcal G}=M_{\mathrm{out}}M_{\mathrm{in}}^{-1}$.
This reconstruction uses the simulated density matrices directly and
does not include measurement sampling or state-preparation errors.

We define the post-gate noise channel by
\begin{equation}
 \mathcal E=\mathcal G\circ\mathcal U_\star^{-1},
 \qquad \mathcal G=\mathcal E\circ\mathcal U_\star,
 \qquad \mathcal U_\star(\rho)=U_\star\rho U_\star^\dagger.
 \label{eq:methods_postgate}
\end{equation}
Thus the noise is placed after the ideal gate. All subsequent error
matrices refer to $\mathcal E$, whose ideal limit is the identity channel.
Using the standard two-qubit Pauli operators
$P_\mu\in\{I,X,Y,Z\}^{\otimes2}$, ordered as
$II,IX,IY,IZ,XI,\ldots,ZZ$, we write
\begin{equation}
 \mathcal E(\rho)=\sum_{\mu,\nu=0}^{15}
       \chi_{\mu\nu}P_\mu\rho P_\nu^\dagger,
 \qquad
 R_{\mu\nu}=\frac14\Tr\!\left[P_\mu^\dagger\mathcal E(P_\nu)\right].
 \label{eq:methods_representations}
\end{equation}
Our normalization is $\Tr(P_\mu^\dagger P_\nu)=4\delta_{\mu\nu}$
and $\Tr\chi=1$ for a trace-preserving channel. Complete positivity and
trace preservation are checked using the Choi matrix
$J(\mathcal E)=\sum_{i,j}\ket i\bra j\otimes\mathcal E(\ket i\bra j)$,
which must satisfy $J\succeq0$ and $\Tr_{\mathrm{out}}J=I_4$.
Trace normalization of $\chi$ alone is not used as a physicality test.

For generator extraction and channel-distance calculations, we project
the reconstructed Choi matrix onto these two constraints in Frobenius
norm using Dykstra's alternating-projection algorithm. The unprojected
matrices are retained to assess the effect of this numerical correction.

\subsection{Error-generator decomposition}
\label{subsec:error_generator}

The effective error generator per gate is represented by the principal
matrix logarithm $K=\log R$. We record the imaginary remainder and the
reconstruction residual $\|e^K-R\|_F$. This finite-gate representation
does not assume that the reduced spin dynamics forms a time-homogeneous
Lindblad semigroup.

To distinguish Hamiltonian and non-Hamiltonian components, we use the
complete elementary-generator basis of Ref.~\cite{Markovian_Errors}.
For nonidentity Pauli operators $P,Q$, the convention is
\begin{align}
 \Ham_P(\rho)&=-\ii[P,\rho], &
 \mathcal S_P(\rho)&=P\rho P-\rho,\\
 \mathcal C_{P,Q}(\rho)&=P\rho Q+Q\rho P
            -\tfrac12\{PQ+QP,\rho\},\\
 \mathcal A_{P,Q}(\rho)&=\ii\left(P\rho Q-Q\rho P
            +\tfrac12\{[P,Q],\rho\}\right).
\end{align}
Their PTM representations span the 240-dimensional space of
trace-preserving, Hermiticity-preserving two-qubit generators. We obtain
the coefficients by linear least squares in
\begin{equation}
 \mathcal K=\sum_{P\ne II}h_P\Ham_P
     +\sum_{P\ne II}s_P\mathcal S_P
     +\sum_{P<Q}\left(c_{P,Q}\mathcal C_{P,Q}
                      +a_{P,Q}\mathcal A_{P,Q}\right).
 \label{eq:methods_full_generator}
\end{equation}
These are signed decomposition coefficients; they are not independently
constrained to be positive Lindblad rates.

We also examine a reduced three-component description,
\begin{equation}
 \mathcal K_3=h_{XX}\Ham_{XX}
           +\gamma_{XX}\Diss[XX]
           +\gamma_{\mathrm{col}}\Diss[S_x].
 \label{eq:methods_reduced_generator}
\end{equation}
The Hamiltonian contribution is obtained from the antisymmetric part of
$K$; the two dissipative coefficients are fitted to the remaining
symmetric part by nonnegative least squares. We report the reconstruction
residual rather than assuming that this reduced model exhausts the
channel. Its fitted $\gamma_{XX}$ and $\gamma_{\mathrm{col}}$ are
distinguished from the coefficients of the complete decomposition.
In particular,
$\Diss[S_x]=\mathcal S_{XI}+\mathcal S_{IX}+\mathcal C_{XI,IX}$,
so a collective $X$ contribution includes a cross term absent from a
model containing only independent $XI$ and $IX$ jumps.

\subsection{Drive-amplitude calibration}
\label{subsec:amplitude_calibration}

For amplitude feedback, the Hamiltonian coefficients are extracted by
least-squares projection of $(K-K^{\mathsf T})/2$ onto the Pauli
Hamiltonian generators. With our target-gate convention, a pure angle
error $U_{\mathrm{actual}}=\exp[\ii(\phi_\star-\epsilon)XX]$
gives $h_{XX}=\epsilon$; positive $h_{XX}$ therefore denotes
underrotation. Guided by the leading-order scaling
$\phi\propto A^2$, we update the rectangular-pulse amplitude as
\begin{equation}
 A_{k+1}=A_k
 \sqrt{\frac{\phi_\star}{\phi_\star-h_{XX}^{(k)}}}.
 \label{eq:methods_calibration}
\end{equation}
The update is used on the branch $\phi_\star-h_{XX}^{(k)}>0$.
After each update, the master equation and all 16 tomography inputs are
recomputed. Consequently, the calibrated channel includes changes in
both its coherent and nonunitary components. The quadratic relation is
an update estimate, not an exact expression for the corrected dynamics.

Calibration is performed independently at each $\bar n$, holding the
detuning, gate duration, and dissipative rates fixed. In particular, the
scattering rates are held fixed in this baseline comparison. The feedback
target is $|h_{XX}|<2\times10^{-3}$, with at most two updates in the
baseline protocol. The achieved residual is retained for every condition.
% AUTHOR CHECK: Keep the protocol aligned with the final figure dataset.
% The current five-point before/after data use one update at each nbar.

\subsection{Fock-state-resolved analysis of thermal errors}
\label{subsec:fock_analysis}

To isolate the Hamiltonian contribution of thermal motion, we set all
explicit dissipative rates to zero and replace the thermal initial state
by individual Fock states. Since $H(t)$ commutes with $S_x$, the unitary
propagator has the conditional form
\begin{equation}
 U(T)=\sum_{s\in\{-2,0,2\}}\Pi_s\otimes U_s(T),
 \qquad
 H_s(t)=\hbar sA(t)
   \left[B e^{-\ii\Phi(t)}+B^\dagger e^{\ii\Phi(t)}\right],
\end{equation}
where $\Pi_s$ projects onto a collective-spin eigenspace.
The $s=0$ branch is stationary in this interaction picture. We calculate
the overlap between the $s=2$ and $s=0$ motional branches,
\begin{equation}
 c_n=\bra n U_2(T)\ket n=r_n e^{2\ii\theta_n}.
 \label{eq:methods_fock_coherence}
\end{equation}
The phase, unwrapped across $n$, defines the effective Fock-conditioned
$XX$ angle $\theta_n$. The magnitude $r_n$ measures the loss of spin
coherence caused by distinguishable conditional motional states. For a
coherent superposition of the two spin branches and a pure Fock input,
$r_n<1$ signals residual spin--motion entanglement.

Using the thermal probabilities in Eq.~\eqref{eq:methods_thermal}, we
evaluate
\begin{align}
 \bar\theta&=\sum_n p_n\theta_n,
 &\sigma_\theta^2&=\sum_n p_n(\theta_n-\bar\theta)^2,\\
 Z&=\sum_n p_n e^{2\ii\theta_n},
 &C&=\sum_n p_n c_n.
\end{align}
The phase-only quantity $Z$ isolates angular dispersion, whereas $C$
also contains residual motional distinguishability. Useful diagnostics
are $h_{XX}^{\mathrm{Fock}}=\phi_\star-\tfrac12\arg C$ and
$\gamma_{\mathrm{phase}}=-\tfrac12\log|Z|$, with the angular branch
chosen continuously near the target. For a narrow angular distribution,
$\gamma_{\mathrm{phase}}\simeq\sigma_\theta^2$.
These diagnostics are compared with the reconstructed generators; a
single branch overlap is not assumed to determine the entire two-qubit
channel. In particular, a collective dissipative interpretation is
assessed using the complete generator and the residual of
Eq.~\eqref{eq:methods_reduced_generator}.

\subsection{Comparison with the Pauli-twirled channel}
\label{subsec:pauli_comparison}

For each channel before and after calibration, we construct its own
Pauli-twirled approximation,
\begin{equation}
 \mathcal T(\mathcal E)(\rho)
 =\frac1{16}\sum_{P\in\{I,X,Y,Z\}^{\otimes2}}
   P^\dagger\mathcal E(P\rho P^\dagger)P
 =\sum_\mu\chi_{\mu\mu}P_\mu\rho P_\mu^\dagger.
 \label{eq:methods_twirl}
\end{equation}
Thus $\chi_{\mathrm P}=\operatorname{diag}(\chi)$ and
$R_{\mathrm P}=\operatorname{diag}(R)$. Twirling is applied to the
channel, not to its logarithmic generator. The average gate infidelity
relative to $U_\star$ is
\begin{equation}
 r=1-F_{\mathrm{avg}}(\mathcal G,\mathcal U_\star)
   =\frac45(1-\chi_{II,II}).
\end{equation}
It is unchanged by Pauli twirling and therefore cannot quantify the
accuracy of that approximation.

We instead evaluate the diamond distance
\begin{equation}
 \epsilon_{\mathrm P}
   =\frac12\|\mathcal E-\mathcal T(\mathcal E)\|_\diamond,
 \label{eq:methods_diamond}
\end{equation}
using primal and dual semidefinite programs~\cite{WatrousDiamond} and
retaining their numerical bounds and gap. As complementary diagnostics,
we report
\begin{equation}
 D_F=\|R-R_{\mathrm P}\|_F,\qquad
 f_F=\frac{D_F}{\|R-I_{16}\|_F},\qquad
 D_J=\frac12\left\|\frac{J(\mathcal E)-J(\mathcal T(\mathcal E))}{4}
        \right\|_1.
 \label{eq:methods_distances}
\end{equation}
Here $\|\cdot\|_1$ is the trace norm and $f_F$ is undefined for the
identity channel. These distances describe the error in replacing a
single physical channel by its Pauli approximation. Logical error rates
of a particular QEC circuit are not computed in this analysis.

\subsection{Numerical settings}
\label{subsec:numerics}

The baseline parameters are listed in Table~\ref{tab:methods_parameters}.
The current before/after channel comparison uses
$\bar n\in\{0.01,1,2,3,4\}$. The equations are integrated with QuTiP,
using absolute and relative solver tolerances of $10^{-12}$ and
$10^{-9}$. Time-dependent Hamiltonian coefficients are sampled at 500
uniformly spaced times over the gate and linearly interpolated; this
coefficient grid is distinct from the adaptive integration steps.
The simulation uses dimensionless values $A_0=0.125$, $\delta=0.5$, and
$T_{\mathrm{sim}}=4\pi$. Physical dissipative rates are multiplied by
$T/T_{\mathrm{sim}}$ when expressed in these simulation units.

\begin{table}[htbp]
 \centering
 \caption{Baseline physical parameters. $A_0$ denotes the effective
 sideband coupling defined in Eq.~\eqref{eq:methods_hamiltonian}.}
 \label{tab:methods_parameters}
 \begin{tabular}{ll}
  \toprule
  Parameter & Value\\
  \midrule
  Gate duration $T$ & $100\,\mu\mathrm{s}$\\
  Detuning $\delta/(2\pi)$ & $10\,\mathrm{kHz}$\\
  Initial coupling $A_0/(2\pi)$ & $2.5\,\mathrm{kHz}$\\
  Lamb--Dicke parameter $\eta$ & $0.1$\\
  Motional diffusion rate $\Gamma_{\mathrm h}$ & $10\,\mathrm{s}^{-1}$\\
  Motional dephasing rate $\Gamma_{\mathrm m}$ & $18\,\mathrm{s}^{-1}$\\
  Spin dephasing parameter $T_2^*$ & $0.3\,\mathrm{s}$\\
  Rayleigh coefficient $\Gamma_{\mathrm{Ray}}$ & $3\,\mathrm{s}^{-1}$\\
  Raman coefficient $\Gamma_{\mathrm{Ram}}$ & $1\,\mathrm{s}^{-1}$\\
  \bottomrule
 \end{tabular}
\end{table}

The oscillator is truncated at a finite dimension $N_{\mathrm{cut}}$,
and the retained thermal probabilities are renormalized. The baseline
QPT implementation chooses this dimension from an occupation-based rule:
with $b=2|A|/\delta$, $\mu=\bar n+b^2$, and
$v=(2\bar n+1)b^2+\bar n(\bar n+1)$, it uses
\begin{equation}
 N_{\mathrm{cut}}=\max\!\left[
 \left\lceil\mu+\bigl(10+2\log(1+\bar n)\bigr)\sqrt v\right\rceil,
 \lfloor5\bar n\rfloor+20,\,2\right]
 \quad(\bar n>0).
\end{equation}
This is a numerical cutoff prescription, not an error bound.
For the Fock-resolved calculation, the largest initial $n$ is selected
so that the omitted thermal probability is below $10^{-5}$, and 24
additional oscillator levels are included above it.
% AUTHOR CHECK: Before submission, add actual cutoff/coefficient-grid/ODE
% convergence results for the final reported observables and conditions.
% Solver tolerances and the cutoff heuristic do not establish convergence.

The CPTP projection uses a tolerance of $10^{-11}$ and at most 5000
iterations, with its constraint residuals recorded. The diamond-distance
calculation uses CVXPY with the SCS solver and tolerance $10^{-8}$.
Projection sensitivity and the SDP gap quantify numerical effects in
the channel post-processing; they do not bound errors due to the
Hamiltonian approximation or oscillator truncation. Finally, the
reconstructed channel assumes an initially factorized spin--motion state
for each gate. Applying it repeatedly in a circuit would require a
separate assessment of motional memory between gates.
% END METHODS

\begin{thebibliography}{9}
\bibitem{QPT}
I.~L. Chuang and M.~A. Nielsen,
``Prescription for experimental determination of the dynamics of a quantum
black box,'' J. Mod. Opt. \textbf{44}, 2455--2467 (1997).
\href{https://doi.org/10.1080/09500349708231894}{doi:10.1080/09500349708231894}.

\bibitem{Markovian_Errors}
R. Blume-Kohout, M.~P. da Silva, E. Nielsen, T. Proctor,
K. Rudinger, M. Sarovar, and K. Young,
``A Taxonomy of Small Markovian Errors,''
PRX Quantum \textbf{3}, 020335 (2022).
\href{https://doi.org/10.1103/PRXQuantum.3.020335}{doi:10.1103/PRXQuantum.3.020335}.

\bibitem{WatrousDiamond}
J. Watrous, ``Simpler semidefinite programs for completely bounded norms,''
\href{https://arxiv.org/abs/1207.5726}{arXiv:1207.5726}.
\end{thebibliography}
\end{document}
